\documentclass[10pt,a4paper]{amsart}
\usepackage[margin=19mm]{geometry}
\usepackage{amsmath,amssymb,mathtools}
\usepackage[scr=rsfs,cal=euler]{mathalfa}
\usepackage{xfrac}
\usepackage[unicode,hidelinks,bookmarks=false]{hyperref}
\newcommand{\ie}{i.e.\ }
\newcommand{\cf}{cf.\ }
\newcommand{\resp}{respectively\ }
\newcommand{\Subsection}{\subsection}
\providecommand{\nobreakdash}{\nobreak\hbox{-}}
\setlength{\emergencystretch}{2em}
\multlinegap0pt
\usepackage[shortlabels]{enumitem}
\usepackage{mathtools}
\usepackage{xcolor}
\usepackage{tikz}
\usepackage{float}
\usepackage{caption}
\newtheorem{lemma}{Lemma}
\title{Decomposition of real numbers into sums of {L\"uroth} sets}
\author{Maiken Gravgaard, Ying Wai Lee}
\date{Source: arXiv:2506.12513v2 (CC BY 4.0)}
\begin{document}
\maketitle

\noindent\emph{Source and licence.} Decomposition of real numbers into sums of {L\"uroth} sets. Version arXiv:2506.12513v2 (CC BY 4.0), distributed by arXiv under the Creative Commons Attribution 4.0 International licence, http://creativecommons.org/licenses/by/4.0/, as recorded in the arXiv licence metadata for this exact version and shown on the abstract page https://arxiv.org/abs/2506.12513v2. Full licence notice: \texttt{licenses/arxiv-2506.12513-CC-BY-4.0.txt}.\par\smallskip

\noindent\textit{Editorial scope.} Counted unit: the article's lemma C (source0.tex lines 1252--1257). The article's complete original proof of it is delivered with it (source0.tex lines 1353--1463, one \texttt{proof} environment of 111 lines, which the article places away from the statement). No mathematics is written, repaired or re-derived: the statement and the proof are the article's own lines, in the article's own notation, and no retained line is substituted or re-flowed.  No further source-local context is needed: every cross-reference of every retained line resolves inside the two delivered spans.  Nothing was omitted from any retained span, so the package carries no omission record.  No retained line contains a citation, so the wrapper carries no bibliography and no bibliography entry.  No retained line contains a figure environment, an image-inclusion command or drawing code, so no image is delivered and the package declares no asset dependency.  Editorial formatting only, in addition: an AMS \texttt{amsart} preamble that restates the article's own declarations and macros used by the retained lines, namely the environments \texttt{lemma} on separate counters, together with the standard packages those lines need.  The retained environments print in this document as Lemma 1 in order of appearance, whereas the article prints the same items with the numbering of its own full document.  The complete native source of the article as distributed in the arXiv e-print for this exact version is bundled unchanged as \texttt{sources/179874/source0.tex}.
	\begin{lemma}\label{lemma: tau C}
		For any $N_1\geq3$, $N_2\geq N_1+2$ and $x\in[0,N_2-N_1-2]$,
		\begin{align*}
			f_0(x)\geq f_1(x).
		\end{align*}
	\end{lemma}

	\begin{proof}[Proof of Lemma~\ref{lemma: tau C}]
		Suppose $N_1\geq 3$ and $N_2-N_1-2\geq0$. For any $x\geq0$,
		\begin{align*}
			D(x)=\frac{c_0+c_1x+c_2x^2+c_3x^3}{C_{N_1,N_2}(N_1+x-1)(N_1+x)(N_1+x+1)},
		\end{align*}
		where the coefficients $c_0,c_1,c_2,c_3$ and the constant $C_{N_1,N_2}>0$ are in terms of $N_1$ and $N_2$,
		\begin{align*}
			c_0&=-{N_1}^5{N_2}+{N_1}^5+{N_1}^4 {N_2}^2-2 {N_1}^4-2 {N_1}^3 {N_2}^2+5 {N_1}^3 {N_2}-{N_1}^3-{N_1}^2 {N_2}+{N_1}^2-{N_1} {N_2}^2\\
			&\qquad-2 {N_1} {N_2}+4 {N_1}+2 {N_2}^2-3 {N_2}-1, \\
			c_1&=-3 {N_1}^4 {N_2}+3 {N_1}^4+2 {N_1}^3 {N_2}^2+{N_1}^3 {N_2}-5 {N_1}^3-3 {N_1}^2 {N_2}^2+8 {N_1}^2 {N_2}-2 {N_1}^2-{N_1} {N_2}^2 \\
			&\qquad-{N_1} {N_2}+3 {N_1}+{N_2}^2-3 {N_2}+1, \\
			c_2&=-3 {N_1}^3 {N_2}+3 {N_1}^3+{N_1}^2 {N_2}^2+2 {N_1}^2 {N_2}-4 {N_1}^2-{N_1} {N_2}^2+4 {N_1} {N_2}-2 {N_1}-{N_2}^2+{N_2}+1, \\
			c_3
			&=-{N_1}^2{N_2}+{N_1}^2+{N_1} {N_2}-{N_1}+{N_2}-1 \\
			&=-({N_1}^2-N_1-1)(N_2-1)<0, \\
			C_{N_1,N_2}
			&= ({N_1}^2-N_1-1)({N_2}^2-N_2-1).
		\end{align*}
		Define $E:\mathbb{R}\to\mathbb{R}$ to be the cubic equation in the numerator of $D$, that is for any $x\in\mathbb{R}$,
		\begin{align*}
			E(x)\coloneqq{c_0+c_1x+c_2x^2+c_3x^3}.
		\end{align*}
		It now suffices to prove the following properties for $E$:
		\begin{enumerate}[A.]
			\item $c_3<0$; hence $\lim_{x\to-\infty} E(x) =+\infty$ and $\lim_{x\to+\infty} E(x) =-\infty$;
			\item $E(-N_1)<0$.
			\item $E(0)>0$;
			\item $E(N_2-N_1-2)>0$;
		\end{enumerate}
		Indeed the above implies that the cubic polynomial $E$ has two negative roots and one positive root that is strictly greater than $N_2-N_1-2$, thus for any $x\in[0,N_2-N_1-1]$, $E(x)>0$ and consequently $D(x)>0$.
		
		\begin{enumerate}[A.]
			\item It is clear from the definition of the coefficient $c_3$ above. 
			
			\item 
			By taking $N_2\geq N_1+2$ we have,
			\begin{align*}
				E(-N_1)
				&=c_0-c_1{N_1}+c_2{N_1}^2-c_3{N_1}^3 \\
				&=\left(-1+3{N_1}-{N_1}^2\right)+\left(-3+{N_1}+{N_1}^2\right){N_2}+\left(2-2{N_1}\right){N_2}^2 \\
				&\leq \left(-1+3{N_1}-{N_1}^2\right)+\left(1-{N_1}-{N_1}^2\right)N_2 \\
				&\leq 1+2{N_1}-4{N_1}^2-{N_1}^3 \\
				&<0.
			\end{align*}
			
			\item 
			$E(0)$ is given by
			\begin{align*}
				E(0)
				&= c_0\\
				&=-{N_1}^5N_2+{N_1}^5+{N_1}^4{N_2}^2-2{N_1}^4-2{N_1}^3{N_2}^2+5{N_1}^3N_2-{N_1}^3 \\
				&\qquad-{N_1}^2N_2+{N_1}^2-N_1{N_2}^2-2N_1N_2+4N_1+2{N_2}^2-3N_2-1 \\
				&=
				{N_2}^2({N_1}^4-2{N_1}^3-N_1+2)+N_2(-{N_1}^5+5{N_1}^3-{N_1}^2-2N_1-3)\\
				&\qquad+({N_1}^5-2{N_1}^4-{N_1}^3+{N_1}^2+4N_1-1).
			\end{align*}
			Suppose $N_1\geq 3$ and $N_2 \geq N_1+2$. the constant term and the coefficients of ${N_2}$ and ${N_2}^2$ are positive, negative and positive respectively. Thus, we have
			\begin{align*}
				c_0 &\geq {N_2}^2\left({N_1}^4-2{N_1}^3-N_1+2\right)+N_2\left(-{N_1}^5+5{N_1}^3-{N_1}^2-2N_1-3\right) \\
				&= N_2\left(N_2\left({N_1}^4-2{N_1}^3-N_1+2\right)+\left(-{N_1}^5+5{N_1}^3-{N_1}^2-2N_1-3\right)\right) \\
				&\geq  (N_1+2)\left({N_1}^4-2{N_1}^3-N_1+2\right)+\left(-{N_1}^5+5{N_1}^3-{N_1}^2-2N_1-3\right) \\
				&= (N_1+1)\left({N_1}^2-3N_1+1\right) \\
				&>0.
			\end{align*}
			
			\item 
			Suppose $N_1\geq3$ and $N_2\geq N_1+2\geq5$. Calculating $E(N_2-N_1-2)$ we get
			{
				\begin{align*}
					E(N_2-N_1-2)
					&=c_0+c_1(N_2-N_1-2)+c_2(N_2-N_1-2)^2+c_3(N_2-N_1-2)^3 \\
					% &=9+13{N_1}-11{N_1}^2+(-16-12 {N_1} +14 {N_1}^2) {N_2}+(8+4 {N_1}-6 {N_1}^2){N_2}^2 +(-1-{N_1}+{N_1}^2){N_2}^3 \\
					&=(9-16{N_2}+8{N_2}^2-{N_2}^3)\\
					&\qquad+(13-12{N_2}+4{N_2}^2-{N_2}^3){N_1}+(-11+14{N_2}-6{N_2}^2+{N_2}^3){N_1}^2
				\end{align*}
			}
			Now we consider the polynomial below for $x\geq3$ and $y\geq5$ 
			\begin{align*}
				P(x, y) = a_y x^2+b_y x+c_y,
			\end{align*}
			where
			\begin{align*}
				a(y) = -11+14y-6y^2+y^3, &&
				b(y) = 13-12y+4y^2-y^3, &&
				c(y) = 9-16y+8y^2-y^3.
			\end{align*}
			Note that $a$ is increasing and that for $y\geq5$,
			\begin{align*}
				a(y)=(y-2)^3+2(y-2)+1\geq34.
			\end{align*}
			Thus, for any fixed $y\geq5$, $P(x, y)$ is quadratic in $x $ and forms an upward parabola. The partial derivative of $P$ with respect to $x$ is given by
			\begin{align*}
				\frac{\partial P}{\partial x}(x,y)&=2a(y)x+b(y).
			\end{align*}
			Evaluating at $x=3$, we see that
			\begin{align*}
				\frac{\partial P}{\partial x}(3,y)
				&=5y^3-32y^2+72y-53 \\
				&=5\left(y-\frac{32}{15}\right)^3+\frac{56}{15}\left(y-\frac{32}{15}\right)+\frac{2369}{675}>0.
			\end{align*}
			This tells us that for fixed $y\geq 5$, $P(x,y)$ is increasing on $x\in[3,\infty)$. Hence, for any $x\geq3$ and $y\geq5$
			\begin{align*}
				P(x,y)
				&\geq P(3,y) \\
				&=(y-3) (5 y^2-19 y+17) \\
				&=(y-3)\left(5 (y-5)^2+31 (y-5)+47\right) \\
				&>0.
			\end{align*}
			
		\end{enumerate}
	\end{proof}

\end{document}
